\documentclass{article}
\title{文章の技術 — 読ませる文章の書き方}
\author{TeX 図書館}

\begin{document}

\section{書く前に決める}

書く技術は、文を直す技術だけでは完成しない。最初に決めることが3つある。目的、読者、読後させたい行動である。ここが曖昧だと、どんなに上手に書いても文章は曖昧なままだ。

\subsection{目的は変化で書く}

目的とは、この文章を書いて何が起きればよいか、である。「情報共有」は目的ではない。共有しても何も変わらないからだ。承認を得る、依頼を通す、誤解を解く。目的は変化で書く。

\begin{quote}
\textbf{【例】}\ 悪い例: 「新人向けに会議のルールを説明する」 / 良い例: 「新人が初回の会議で、自分の案件を1つ発言できるようにする」
\end{quote}

\begin{quote}
\textbf{【注意】}\ 「共有します」「検討します」は目的の言葉に見えて、実際には状態である。読んだあとに何が変わっているかを書けないなら、それは目的ではなく、書き手の都合である。
\end{quote}

\subsection{読者を3つの質問で絞る}

読者は毎回違う。3つだけ質問すれば絞れる。誰が読むか。その人は何をすでに知っているか。その人は何を気にしているか。この3つで、載せる情報も語る順番も決まる。

\begin{quote}
\textbf{【例】}\ 同じ問い合わせ対応の報告でも、読者で載せるものが変わる。上司向けなら、影響範囲と再発防止を先に置く。顧客向けなら、復旧時刻と今後の対策を先に置く。書き手が知りたいことではなく、読者が決めたいことが、載せ物を選ぶ。
\end{quote}

\subsection{読後させたい行動を最初に置く}

読者は最後まで読まない。最初の2〜3文で「読む価値がある」と判断する。だから、読後に期待する行動を最初の段落に置く。奥に埋めた依頼は、たいてい届かない。

\begin{quote}
\textbf{【例】}\ 悪い例: 「先日、件Aの経緯について社内で検討いたしましたので、ご報告申し上げます。」 / 良い例: 「件Aは、木曜の正午までにA案とB案のどちらかを選ぶ必要があります。判断材料を3点まとめました。」
\end{quote}

\begin{quote}
\textbf{【演習】}\ 「部署に新しい議事録の様式を導入したい」という思いを、目的・読者・読後させたい行動の3点で書き直せ。
\end{quote}

\textbf{解答の指針:} 目的は「毎週の決定事項が、口頭確認なしで残るようにする」のように変化で書く。読者は「議事録を書く担当者。作業時間が減ることを気にする」まで具体化する。行動は「来週から新しい様式で1件、書き切ること」のように書き、最初の段落に置ける一文にする。

\section{構成}

決めた材料は、順番を変えるだけで読みやすさが変わる。構成の基本は逆三角形である。結論から書き、根拠を重ね、詳細は最後に置く。読者は結論を探しており、途中の物語を読む余裕はない。

\subsection{逆三角形で結論を先に出す}

\begin{quote}
\textbf{【例】}\ 悪い例: 「先月、A社との取引条件を見直すことになりました。背景には、原材料価格の変動があります。その詳細は次の通りです。結論として、単価を5\%見直したいと考えています。」 / 良い例: 「A社との単価を5\%見直したいです。理由は原材料価格の変動で、先月の調達費は前年比で12\%上がりました。詳細を3点示します。」
\end{quote}

内容は同じなのに、読む負荷がまるで違う。悪い例では読者は3段落目まで結論を待つ。待たされた読者は、その間に閲覧をやめる。

\subsection{リードで読む理由を渡す}

最初の2〜3文の仕事は3つある。読む理由を与えること。読者との接点を示すこと。この先の見取り図を渡すこと。この3つが終われば、リードとしては完成である。

\begin{quote}
\textbf{【注意】}\ 前置きで埋めない。「ご担当者様。いつもお世話になっております。平素より格別のご高配を賜り、厚く御礼申し上げます。本件につきましては」まで書いても、用件は1文字も進んでいない。挨拶は1行で足りる。用件は2行目から始める。
\end{quote}

\subsection{段落は1つ1論点}

段落の最初の一文は、その段落の要約である。最後の一文は、次の段落への橋になる。2つの論点を1段落に入れると、要約の一文が書けない。書けなければ、分割すべき段落である。

\begin{quote}
\textbf{【例】}\ 悪い例: 「新機能は操作が簡単で、問い合わせも減った。一方で開発コストがかかり、来期の予算に影響する。」 / 良い例: 「新機能のおかげで、問い合わせは月20件から7件に減った。一方で開発コストは300万円かかった。予算への影響を次の段落で示す。」
\end{quote}

\subsection{見出しは約束である}

見出しは、次に読む段落の内容を言い当てること。これが約束である。守らないと、読者は見出しの直後で迷い、読み直しを始める。「はじめに」「いろいろ」「補足」のような、内容を言わない見出しは約束になっていない。名詞句で、中身を言い切る。

\section{文を書く}

構成が決まれば、次は一文ずつである。文の技術は4つに集約される。一文一義、主語と述語の距離、体言止めの使い所、受身形と敬語の重さだ。

\subsection{一文に判断は1つ}

一文は一つの判断だけを持つ。接続詞が2つ並んだら、分割のサインである。

\begin{quote}
\textbf{【例】}\ 悪い例: 「A案は費用が安いが導入に3週間かかるため、B案は速いが高価なので、両方の利点を持つC案を検討すべきと考える。」 / 良い例: 「A案は安いが、導入に3週間かかる。B案は速いが、高価だ。C案なら両方の利点を持つ。検討したい。」
\end{quote}

短くした結果、文は4つになった。しかし読者の理解は速い。1文の中で3回も判断を迫られるより、判断を1つずつ渡すほうが、負荷が低いからである。

\subsection{主語と述語を近づける}

修飾語は、修飾される語のすぐそばに置く。主語が遠い文では、読者が述語を探す間に息が尽きる。さらに主語を落とすと、誰がやったのかが消える。

\begin{quote}
\textbf{【例】}\ 悪い例: 「確認した資料を、添付して送信した。」 / 良い例: 「部長が確認した資料を、私は添付して送信した。」
\end{quote}

悪い例では、誰が確認したのか一読では決まらない。書き手には分かっていることが、読み手には分かっていない。主語の省略は、文を短くする技術であり、誤解を生む技術でもある。

\subsection{体言止めは調味料}

体言止めは効く。強調と切れ、そしてリズムが生まれる。ただし連発すると宣伝文句になる。地の文で3連続したら、1つは言い切りに直す。

\begin{quote}
\textbf{【注意】}\ 「新しい機能。使いやすい画面。無料の試用。」は広告のキャッチコピーにはいい。説明文には合わない。体言止めは力点であり、通し続ければ読者は疲れてしまう。
\end{quote}

\subsection{受身形と敬語は重い}

「対応が検討されています」は、誰が検討しているかを消している。責任の所在を消す文は、安心させようとして不安を生む。敬語は相手への礼であり、情報の削りではない。

\begin{quote}
\textbf{【例】}\ 悪い例: 「ご対応につきましては、社内で検討が重ねられております。」 / 良い例: 「対応は、来週金曜までに可否をお答えします。判断者は営業部の山田です。」
\end{quote}

\begin{quote}
\textbf{【演習】}\ 「A案は安いが手間がかかるので、B案も含めて比較した上で、来週の会議で決めるべきだと思う。」を、一文一義に書き直せ。
\end{quote}

\textbf{解答の指針:} 判断を3つに割る。「A案は安いが手間がかかる。」「B案も含めて比較したい。」「来週の会議で決める。」接続詞の数が減った文ほど、読者の理解は速い。末尾の「と思う」は、意見として残すか落とすかを自分で決める。

\section{リズム}

同じ内容でも、文の長さが並ぶと読者は眠る。リズムは技巧ではなく、集中を保つための装置である。作り方は4つある。

\subsection{長短を交互に置く}

長い文の次に短い文を置く。短い文が着地点になる。

\begin{quote}
\textbf{【例】}\ 「この提案は、来期の予算編成、採用計画、拠点の統合という3つの大きな判断に影響する。先送りはできない。」長い文で重みを作り、短い文で着地させる。逆にすると、締めがぼやける。
\end{quote}

\subsection{読点は呼吸である}

読点は文法ではなく、呼吸の印だ。息が尽きる前に切る。切る場所に迷ったら、文が長すぎる。

\begin{quote}
\textbf{【例】}\ 悪い例: 「昨日駅前に寄った時に閉店セールをしていた店で買ったパンはあまりに美味しくて明日も買いに行きたいくらいだった。」 / 良い例: 「昨日、駅前の店でパンを買った。閉店セールだった。あまりに美味しくて、明日も買いに行きたい。」
\end{quote}

\subsection{同じ語尾の3連続を避ける}

「〜である。〜である。〜である。」は単調だ。同じ語尾が3つ並んだら、文を合体させるか、疑問形か言い切りに1つ変える。

\begin{quote}
\textbf{【注意】}\ 文書の中で「だ/である」と「です/ます」を混ぜるのは、たしかに乱れである。一方、語尾の変化(〜た。〜ている。〜か)はリズムの範囲だ。混在と変化は別物である。文体は統一し、語尾は変える。
\end{quote}

\subsection{段落は3〜4文で切る}

段落が長いと、読者は途中で戻れなくなる。3〜4文で切る。切り替えの気休めにもなる。速く読む場面では、一画面に一論点が最も速い。

\section{語彙}

言葉の選び方は、翻訳の受け渡しに似ている。抽象語は読み手に翻訳を任せる。具体語は翻訳済みである。迷ったら、具体的なほうを選ぶ。

\subsection{具体語で書く}

\begin{quote}
\textbf{【例】}\ 悪い例: 「本施策により、業務効率が向上します。」 / 良い例: 「毎月の締め作業が、3日から1日に短くなります。」
\end{quote}

「効率が向上する」は、読者が自分で何がどう変わるかを思い浮かべる必要がある。思い浮かべられない読者は、要らないと判断する。数字は、思い浮かべる手間を肩代わりする。

\subsection{動詞で書く}

日本語は、名詞で止まりやすい言語である。「検討を行う」は「検討する」で足りる。「改善が図られる」は「改善する」で足りる。「〜することができます」は「〜できます」である。

\begin{quote}
\textbf{【例】}\ 悪い例: 「新機能の利用が可能となります。」 / 良い例: 「新機能を使えます。」
\end{quote}

名詞で止める文は息が止まる。動詞で終わる文は前に進む。手が動く文章ほど、読者も動きやすい。

\subsection{くどい表現を削る}

削れる言葉は、ほぼ決まっている。「〜ということになる」「〜と考えられる」「〜のような形で」「〜という点については」。これらは意味を足さない。

\begin{quote}
\textbf{【例】}\ 悪い例: 「不具合が修正されたということになりました。」 / 良い例: 「不具合を直しました。」
\end{quote}

\subsection{和語・漢語・外来語のバランス}

和語は身近で、漢語は正確で、外来語は新しい事物に合う。迷ったら和語である。「フォーマット」は多くの場合「様式」で足りる。ただし、その職場で定着した語や、固有の製品名はそのまま使う。

\begin{quote}
\textbf{【注意】}\ 翻訳調に注意する。「我々は、彼の提案に同意する所存でございます」は、日本語としては不自然である。「私たちは、彼の提案に賛成します」で同じ意味が伝わる。外国語の語順をそのまま日本語に落とさない。
\end{quote}

\begin{quote}
\textbf{【演習】}\ 「本件につきましては、迅速な対応が図られる可能性があるということです」を、書き直せ。
\end{quote}

\textbf{解答の指針:} 名詞化(対応が図られる)、くどい表現(ということです)、逃げの条件(可能性)を削る。「できるだけ早く対応します。木曜までにお返事します」のように、誰がいつやるかを動詞で書く。削って意味が変わらなかった言葉は、元々いなかったのである。

\section{推敲}

初稿は下書きである。書いた直後の文章は、まだ読者の手に渡っていない。推敲は3つの作業に分かれる。削る、並べ替える、音読する。順番もこの通りである。

\subsection{まず削る}

目標は、1割削ることである。削る候補は最初に見つかる。「非常に」「大変」「〜ではないかと考えられます」「〜することができる」。これらを消して、意味が変わらなかったら、それはなかった言葉だ。

\begin{quote}
\textbf{【例】}\ 悪い例: 「本日は大変お忙しいところ恐れ入りますが、以下の点についてご確認いただけますと非常に幸いです。」 / 良い例: 「以下の2点をご確認ください。木曜までにお願いします。」
\end{quote}

\subsection{次に並べ替える}

詰まる箇所は、たいてい文の質ではなく配置の問題である。段落の順を入れ替えるだけで、結論が強くなることがある。背景→詳細→結論を、結論→背景→詳細に並べ替えるだけで、同じ文でも読む速度が変わる。

\begin{quote}
\textbf{【例】}\ 悪い例: 詳細な調査結果を5段落並べて、最後に「よって、Xを推奨します」。 / 良い例: 「Xを推奨します」を最初に置き、5段落を根拠として後に置く。文は1つも足さずに、説得力が増す。
\end{quote}

\subsection{音読する}

耳は目より正直である。つっかえる所、息が尽きる所、同じ音が続く所が、すべて弱点だ。声に出したくなかった文は、読者も読みたくない。

\begin{quote}
\textbf{【注意】}\ 画面で流し読みできる文も、声に出すと引っかかることがある。引っかかったら、文を切るか、短い言い切りに変える。音読は、推敲の中で唯一、他人の代わりに試してくれる作業である。
\end{quote}

\subsection{自己診断のチェック項目}

推敲の終わりに、次の項目を上から順に確かめる。1つでも引っかかれば、その場で直す。

\begin{itemize}
\item 1行目だけで、用件と依頼が分かるか
\item 一文に、判断は1つ入っているか
\item 主語と述語が、すぐ結びついているか
\item 同じ語尾が3つ以上、連続していないか
\item 抽象語に、具体や数字が伴っているか
\item 削れる言葉が、まだ残っていないか
\item 事実と意見が、混ざっていないか
\item 最後に、読者への行動が書いているか
\end{itemize}

\section{読ませる工夫}

技術を身につけたあとで、初めて読ませる工夫を考える。読ませる工夫は誇張ではない。事実を、読者が追いやすい順に置くことである。

\subsection{フックは最初の一文}

読む理由を、1文で渡す。数字、断言、問い、意外性。この4つが使える。

\begin{quote}
\textbf{【例】}\ 悪い例: 「今回は時間管理について書きます。」 / 良い例: 「会議は毎週、90分ずつ無駄になっている。ここに1つの対策がある。」
\end{quote}

前者は書き手の予告である。後者は読者への差し入れだ。フックは、読者が得をする予告でなければならない。

\subsection{例えは1つで足りる}

抽象を、読者がすでに知っているものに置き換える。新製品の説明なら、読者の日常の1場面に置く。たとえば「設定が3ステップで終わる」なら、朝のコーヒーを淹れる時間より短い、と言えばよい。

\begin{quote}
\textbf{【注意】}\ 例えがズレると、むしろ混乱する。相手の職種や生活が分からないまま、自分の世界の例えを持っていくのは危険だ。例えは1段落に1つで足りる。2つ目からは、説明になってしまう。
\end{quote}

\subsection{数字と場面は正確に}

「多くの顧客」は「問い合わせの7割」に、「早く対応」は「平均12分で折り返し」に置き換える。数字は正確であること。曖昧な大げさな数字より、控えめで正しい数字のほうが信頼される。場面も1つ描く。誰が、どこで、どう困っていたか。1場面が、3つの形容詞に勝る。

\subsection{抑揚は事実の順で作る}

誇張はしない。問題、転機、解決。この順番が、そのまま抑揚になる。

\begin{quote}
\textbf{【例】}\ 「毎晩、手作業で3時間かかっていた。自動化を試した。今は10分で終わる。」3文だけで抑揚が生まれる。形容詞で盛る必要はない。事実の並びが強い。
\end{quote}

\section{形式別の型と長さ}

型は制約ではない。速さの装置である。場面ごとの型を知っておくと、推敲の手間が減る。ここでは4つの形式を取り上げる。

\subsection{メール}

メールの型は4つである。件名で用件、1行目で結論、依頼は1つ、期限は日付と時刻で書く。

\begin{quote}
\textbf{【例】}\ 悪い例: 件名「ご相談」、本文「お世話になっております。先日の件、いかがでしょうか。進捗おりましたらご連絡いただけますと幸いです。」 / 良い例: 件名「【承認依頼】A社契約書 3/15 17時まで」、本文「A社の契約書、木曜17時までにご確認ください。変更点は価格と納期の2箇所です(赤字)。ご返信がなければ金曜に送付します。」
\end{quote}

社内メールは500字あれば足りる。依頼は冒頭の1段落に集める。用件を文末に置いたメールは、読み飛ばされる。

\subsection{ドキュメント}

社内ドキュメントは、誰が読むか分からないまま書かれる。だから、前提を最初に書く。略語、社内通称、日付の基準。読者が詰まる箇所を先に埋める。章立ては1画面1論点で、見出しの約束を守る。

\begin{quote}
\textbf{【例】}\ 悪い例: 「はじめに 当社の当方針につきましては、従前より大きく変わるものではありません。」 / 良い例: 「この文書は、7月から始まる在宅勤務の規則を説明する。対象は全社員。例外は経営会議の承認のみ。」
\end{quote}

\subsection{SNSと広告文}

1文目がすべてである。スクロールする読者に、残る理由を1行で渡す。抽象は禁止。数字、固有の場面、読者の得を先に置く。1投稿1メッセージで、140字以内。ハッシュタグは最後に1つ。

\begin{quote}
\textbf{【例】}\ 悪い例: 「新しい機能をリリースしました!ぜひお試しください。」 / 良い例: 「ログイン設定が3ステップに短縮しました。今なら、設定は2分で終わります。」
\end{quote}

\subsection{報告書}

報告書の順番は、結論、根拠、詳細の3段である。事実と意見を分けて書く。「思われます」「かもしれません」で逃げない。事実は数値で、判断は担当名で書く。結論は3行以内に収め、詳細は後半に置く。読む人が最後まで読まない前提で、最初の3行にすべての要点を入れる。

\begin{quote}
\textbf{【演習】}\ 件名「ご連絡」、本文「お世話になっております。先日の資料の件で、一部修正が生じましたのでご報告いたします。詳細は添付のとおりです。」を、用件が先になるように書き直せ。
\end{quote}

\textbf{解答の指針:} 件名は「何が・どれが・いつ」を詰める。「【修正報告】提案書 v2 本日17時更新」のようにする。1行目は結論・変更点・読者が取る行動の3点を置く。「提案書を修正しました。変更点は価格と納期の2つです。17時以降の版を使ってください」が目安である。

\section{総合 — 自分の文章に戻る}

技術は、自分の文章に戻って初めて技術になる。この章では、チェックリストと習慣をまとめる。

\subsection{送信前チェックリスト}

下書きを送る前に、次の項目を上から順に確かめる。全部で8項目、所要は2分である。

\begin{itemize}
\item 目的を、変化で書けているか(読んだ後に何が変わるか)
\item 読者を3つの質問で絞ったか(誰か、何を知っているか、何を気にするか)
\item 読後させたい行動を、最初の段落に置いたか
\item 結論は先に出てきているか(逆三角形になっているか)
\item 一文に判断は1つずつか。主語と述語は近くにあるか
\item 同じ語尾が3連続していないか。長短が交互に入っているか
\item 抽象語は具体語に置き換わっているか。動詞で書いているか
\item 削れる言葉を削ったか。声に出して読み返したか
\end{itemize}

\subsection{続けるための習慣}

3つだけを続けるとよい。寝かせる、比べる、写す。

\begin{quote}
\textbf{【例】}\ 悪い例: 書いたその場で、そのまま送る。 / 良い例: 5分置いてから、1回だけ読み返す。直した版と初稿を並べて、どこが変わったかを見る。
\end{quote}

比べる作業は特に効く。自分の初稿と直した版の差が見えれば、「何を削ったか」が身体に入り、次は最初から削って書けるようになる。写すのは、自分の好きな文章である。手で写して声に出すと、その人のリズムが体に入り込む。上手な文章は、暗唱できる文章から始まることも多い。

\begin{quote}
\textbf{【注意】}\ 完璧な一文より、正確で速い一文である。締切のある文章では、推敲する時間そのものを最初に確保する。書く時間3、直す時間1の割合で配分する。直さない文章は、いつまでも下書きのままだ。
\end{quote}

\begin{quote}
\textbf{【総合演習】}
\begin{enumerate}
\item 自分が最近書いたメールを1通思い出し、送信前チェックリストの8項目で自己採点せよ。
\item 「業務効率が向上しました」という一文を、具体語と数字で書き直せ。
\item 結論が文末に来てしまっている報告の冒頭3文を、逆三角形に並べ替えよ。
\end{enumerate}
\end{quote}

\textbf{解答の指針:} 1. 8項目のうち埋まっていたものを数え、欠けた項目を次回の下書きで先に決める。2.「締め作業が3日から1日に」「問い合わせが月20件から7件に」など、読者が思い浮かべられる数字を伴わせる。3. 結論1文を冒頭に移し、根拠と詳細を後ろに並べ替える。文を1つも足さずに、読む速度が上がっているはずである。

\end{document}
